% TOP_PROBLEM: {"id": 30007427, "title": "Almost everywhere Fourier convergence for L log L", "category_id": 8, "category": {"id": 8, "name": "analysis", "display_name": "Analysis", "description": "Limits, continuity, calculus, and function theory.", "slug": "analysis", "order_index": 8, "created_at": "2026-07-31T15:26:25.671Z"}, "status": "open"}
% FIRST_POSTED: 2026-10-06
% ATTEMPT_STATUS: solved
% OPENAI_CLAIM: {"schema_version":1,"source_repo":"https://github.com/openai/math","source_commit":"adc7f1241b42e322a6451854ab7e4b4c146bf78a","release_date":"2026-10-06","match":"full","resolution":"proved","independently_reviewed":false,"second_pass":"agreed","families":[{"family":"075","title":"The $`L\\log L`$ Fourier-convergence conjecture","subject":"Real and complex analysis","manuscripts":[{"title":"Almost-everywhere Fourier convergence in Llog L","pdf_path":"preprints/Almost-everywhere-Fourier-convergence-in-L-log-L-September-23-2026/paper.pdf","readme_path":"preprints/Almost-everywhere-Fourier-convergence-in-L-log-L-September-23-2026/README.md","date":"2026-09-23","theorem_ref":"thm:main (sections/01-introduction.tex:17)","statement_sources":[{"path":"preprints/Almost-everywhere-Fourier-convergence-in-L-log-L-September-23-2026/build/sections/01-introduction.tex","line":17,"label":"thm:main"}]}],"lean":null,"reasoning_trace":null}]}
\documentclass[11pt]{article}
\usepackage[margin=25mm]{geometry}
\usepackage{fontspec}
\setmainfont{Latin Modern Roman}
\usepackage{amsmath,amssymb,mathtools}
\usepackage{unicode-math}
\setmathfont{Latin Modern Math}
\usepackage{enumitem}
\usepackage{xurl,hyperref,bookmark}
\hypersetup{colorlinks=true,urlcolor=blue}
\setcounter{secnumdepth}{0}
\setlength{\parindent}{0pt}
\setlength{\parskip}{5pt}
\setlength{\emergencystretch}{3em}
\begin{document}
\section*{Almost everywhere Fourier convergence for L log L}

\subsection{OpenAI's claim}
\textit{Fourier convergence at the Orlicz endpoint}: For a complex function on the circle satisfying the $L\log L$ integrability condition, the paper claims convergence of its symmetric Fourier sums almost everywhere. The limit is asserted for the ordinary full sequence of partial sums.

\subsection{Mathematical statement of the claimed result}
\paragraph{Mathematical claim 1.} On $\mathbb T=\mathbb R/(2\pi\mathbb Z)$ with normalized measure $\mu$, suppose $\int|f|\log(2+|f|)\,d\mu<\infty$. Set $\widehat f(k)=\int f(x)e^{-ikx}\,d\mu(x)$ and $S_Nf=\sum_{k=-N}^N\widehat f(k)e^{ikx}$. The claimed conclusion is \[S_Nf(x)\longrightarrow f(x)\quad\text{for almost every }x,\quad N=0,1,2,\ldots.\] It applies to complex-valued inputs and the full sequence of symmetric partial sums, under the stronger $L\log L$ hypothesis rather than merely $L^1$.

\subsection{Outline of the argument}
\textit{Fourier convergence at the Orlicz endpoint}: A bounded nonnegative density is first controlled on the set where its spatial maximal average is moderate. The corresponding modulated dyadic sums are stopped on a finite binary tree, giving a good-set integral estimate uniform in the depth. Averaging delayed Haar kernels recovers the Hilbert kernel; doubly exponential height layers then make both the integral costs and exceptional-set measures summable.

\subsection{Scope relative to this problem}
\textit{Fourier convergence at the Orlicz endpoint}: The primary statement says: Every complex L log L function on the circle has its entire sequence of symmetric Fourier partial sums converge to the function almost everywhere. The target retains these quantified objects and hypotheses, so its positive assertion is covered.

With the hypotheses and conventions above, the claimed result addresses the whole mathematical assertion on this page.

\subsection{Formal verification}
\textit{Fourier convergence at the Orlicz endpoint}: The pinned catalogue lists no Lean formalization for this result.

This site has not run the Lean development. These links describe OpenAI’s selected formal statements and should be read with their stated scope.

\subsection{Status caveat}
OpenAI reports that an internal model produced the manuscripts. Their arguments have not been independently verified as proofs by this site. Any portion outside the documented formal scope may contain mistakes, and the catalogue label reports the announced mathematical claim.

\subsection{Sources}
\begin{itemize}
\item OpenAI, \emph{Almost-everywhere Fourier convergence in Llog L}, 2026-09-23; citation key \texttt{OAI:Almost-everywhere-Fourier-convergence-in-L-log-L-September-23-2026}. \href{https://github.com/openai/math/blob/adc7f1241b42e322a6451854ab7e4b4c146bf78a/preprints/Almost-everywhere-Fourier-convergence-in-L-log-L-September-23-2026/paper.pdf}{PDF (GitHub)}; \href{https://github.com/openai/math/raw/adc7f1241b42e322a6451854ab7e4b4c146bf78a/preprints/Almost-everywhere-Fourier-convergence-in-L-log-L-September-23-2026/paper.pdf}{PDF (direct)}; \href{https://github.com/openai/math/blob/adc7f1241b42e322a6451854ab7e4b4c146bf78a/preprints/Almost-everywhere-Fourier-convergence-in-L-log-L-September-23-2026/README.md}{Manuscript page and citation}; \href{https://github.com/openai/math/blob/main/preprints/Almost-everywhere-Fourier-convergence-in-L-log-L-September-23-2026/paper.pdf}{Latest version}; Statement source: \href{https://github.com/openai/math/blob/adc7f1241b42e322a6451854ab7e4b4c146bf78a/preprints/Almost-everywhere-Fourier-convergence-in-L-log-L-September-23-2026/build/sections/01-introduction.tex\#L17}{thm:main}.
\item \href{https://github.com/openai/math/blob/adc7f1241b42e322a6451854ab7e4b4c146bf78a/CONTENTS.md}{OpenAI catalogue entry}.
\end{itemize}
\end{document}
